% ===================== 六、模型检验与分析 =====================
% 结构：引文 -> 二级标题分类 -> 文字 -> longtable -> 分析。
% 子节：误差分析（优先 5 折交叉验证）+ 灵敏度分析（关键参数 ±20% 扰动）。
% 灵敏度数据在求解阶段前置产出，此处直接引用。
\section{模型检验与分析}

为验证本文构建的模型的有效性和稳健性，进行以下检验。

\subsection{误差分析}

\begin{longtable}{>{\centering\arraybackslash}p{0.22\textwidth}>{\centering\arraybackslash}p{0.22\textwidth}>{\centering\arraybackslash}p{0.22\textwidth}>{\centering\arraybackslash}p{0.22\textwidth}}
  \caption{模型误差分析结果}
  \label{tab:error} \\
  \toprule
  指标 & 训练集 & 测试集 & 说明 \\
  \midrule
  \endfirsthead
  \caption{模型误差分析结果（续）} \\
  \toprule
  指标 & 训练集 & 测试集 & 说明 \\
  \midrule
  \endhead
  \bottomrule
  \endfoot
  $R^2$ & ... & ... & 决定系数 \\
  RMSE & ... & ... & 均方根误差 \\
  MAE & ... & ... & 平均绝对误差 \\
  MAPE(\%) & ... & ... & 平均绝对百分比误差 \\
\end{longtable}

如表\ref{tab:error}所示，...

\subsection{灵敏度分析}

\begin{longtable}{>{\centering\arraybackslash}p{0.30\textwidth}>{\centering\arraybackslash}p{0.30\textwidth}>{\centering\arraybackslash}p{0.30\textwidth}}
  \caption{灵敏度分析结果}
  \label{tab:sensitivity} \\
  \toprule
  参数变化 & 输出变化 & 灵敏度评价 \\
  \midrule
  \endfirsthead
  \caption{灵敏度分析结果（续）} \\
  \toprule
  参数变化 & 输出变化 & 灵敏度评价 \\
  \midrule
  \endhead
  \bottomrule
  \endfoot
  ... & ... & ... \\
\end{longtable}

如表\ref{tab:sensitivity}所示，...
